%% Interstitial: Paper 2 -> Paper 3
%% Narrative: the dynamics are characterised abstractly;
%% now we ask how to instantiate and compute them.

Chapter~\ref{chap:paper-two} established that the predictive
operator family $\{K_t\}$ forms a semigroup and is dual to the
Koopman operator.  This is an abstract characterisation: it holds for
any observable dynamics satisfying the program's hypotheses, without
specifying the form of the queries or the structure of the dynamics.

For the program to make contact with data --- with actual time series,
actual measurements, actual algorithms --- it needs a concrete
instantiation.  What queries arise in practice?  What does the
semigroup $\{K_t\}$ look like when the outcome spaces have geometric
structure?  And how can $K_t$ be approximated from finite data?

\medskip

Chapter~\ref{chap:paper-three} provides the instantiation via
\emph{delay query systems}.  The delay query at lag $\tau$ and
window $m$ measures the system by recording $m$ consecutive
observations spaced $\tau$ apart: $Q_{m,\tau}(\omega) =
(\omega(0), \omega(\tau), \ldots, \omega((m-1)\tau))$.  These
queries are natural --- they formalise the Takens embedding strategy
used throughout applied dynamical systems --- and their refinement
structure is governed by the program's directedness conditions.

The main results of Chapter~\ref{chap:paper-three} are:
the delay query system is upper-directed (so the extension theorem
applies), the probabilistic sufficiency criterion identifies when a
delay window is long enough to capture the full predictive content,
and the Markov order $m(\tau, P)$ measures the minimal window length
required.  The operator $K_t$ is approximated by $K_t^N$ computed
from $N$ data points via local linear regression on the delay
embedding.

\medskip

This is where the program meets the world.  The abstractions of the
preceding chapters --- queries, compatibility, predictive kernels,
semigroups --- become an algorithm: observe the system through a delay
window, estimate the predictive operator, track how it evolves.

The passage from Chapter~\ref{chap:paper-two} to
Chapter~\ref{chap:paper-three} is the passage from \emph{dynamics} to
\emph{computation}: from the abstract semigroup structure to the
concrete instantiation that makes it estimable from data.
